% Generated draft: requires mathematical and source review.
\documentclass[11pt]{article}
\usepackage[margin=25mm]{geometry}
\usepackage{fontspec}
\setmainfont{Latin Modern Roman}
\usepackage{amsmath,amssymb,mathtools,unicode-math}
\setmathfont{Latin Modern Math}
\usepackage{enumitem,needspace,xurl,hyperref}
\setlength{\parindent}{0pt}
\setlength{\parskip}{5pt}
\setlength{\emergencystretch}{3em}
\newcommand{\N}{\mathbb{N}}
\newcommand{\Z}{\mathbb{Z}}
\newcommand{\Q}{\mathbb{Q}}
\newcommand{\R}{\mathbb{R}}
\newcommand{\C}{\mathbb{C}}
\newcommand{\F}{\mathbb{F}}
\newcommand{\A}{\mathbb{A}}
\newcommand{\PP}{\mathbb{P}}
\newcommand{\E}{\mathbb{E}}
\newcommand{\eps}{\varepsilon}
\newcommand{\dd}{\,\mathrm{d}}
\begin{document}
\section*{UnsolvedMath ID: 30004792}
\textbf{Source problem code:} catalogue entry
\subsection{Definitions and mathematical statement}
The supplied record proposes an unramified automorphic--spectral correspondence. Here \(q\) denotes the cardinality of the finite field \(\mathbb F_q\). The notation \(\operatorname{Bun}_G\) denotes a moduli stack of principal \(G\)-bundles on an implicit underlying curve, and \(\widehat G\) denotes the intended Langlands dual group. The record does not specify that curve, its hypotheses, the precise class and field of definition of \(G\), or the coefficient theory used to form the dual group and the spaces below. These are unresolved parameters, not additional assumptions supplied by this draft.

The proposed statement is the existence of an isomorphism
\[
\Phi:\operatorname{Func}\bigl(\operatorname{Bun}_G(\mathbb F_q)\bigr)
\xrightarrow{\ \sim\ }
\Gamma\bigl(\operatorname{LocSys}^{\operatorname{arithm}}_{\widehat G},\omega\bigr)
\]
compatible with Hecke operators. The left side is intended to be a space of unramified automorphic functions on isomorphism classes of \(\mathbb F_q\)-points of the bundle stack. The notation \(\operatorname{LocSys}^{\operatorname{arithm}}_{\widehat G}\) is intended to describe a moduli object of arithmetic \(\widehat G\)-local systems on the implicit curve; \(\omega\) denotes its intended dualizing object, and \(\Gamma\) denotes global sections. Hecke compatibility means that \(\Phi\) intertwines the corresponding automorphic and spectral Hecke actions, once those actions have been specified.

The record leaves unspecified the coefficient field, support or growth conditions on functions, the meaning of arithmetic local systems, the geometric or derived structure of their moduli object, and whether the displayed global sections are a complex or a specified cohomology group. It also leaves unspecified the Hecke operators and their spectral action. Consequently, the display records the supplied conjectural target but does not yet define an isomorphism between fully specified mathematical objects.

The record additionally asserts an equivalent formulation: the relevant Frobenius traces of nilpotent \(\operatorname{IndCoh}\) recover the dualizing global sections. Here \(\operatorname{IndCoh}\) refers to an ind-coherent sheaf category, with a nilpotent support condition intended. The underlying stack, precise support condition, Frobenius action, trace construction, and comparison map are not supplied. This alternative formulation and its asserted equivalence therefore remain unresolved parts of the source statement.
\subsection{Short English statement}
The supplied conjecture asks for a Hecke-compatible identification of unramified automorphic functions with global sections of a dualizing object on arithmetic Langlands-dual local systems. The curve, group hypotheses, coefficients, function space, spectral moduli object, and categorical Frobenius-trace formulation need specification before this becomes a precise conjecture.
\subsection{Sources}
\begin{itemize}
\item Local UnsolvedMath registry, permanent ID \(30004792\), title: \emph{Automorphic Functions as Sections on Arithmetic Local Systems}; source problem code \texttt{catalogue entry}, in the Oberwolfach Reports Open Problems collection. No individual proposer is supplied; the proposed year is recorded as 2021.
\item Supplied report citation: \emph{Automorphic Forms, Geometry and Arithmetic (2022)}, Oberwolfach Reports, DOI 10.4171/owr/2021/39, pp.~2089--2156. Source URL and literature-source link labelled ``original Oberwolfach report'': \url{https://doi.org/10.4171/owr/2021/39}. The supplied background describes these reports as documenting mathematical workshops and their open-problem discussions. No specific page or contribution locating this conjecture is supplied.
\item Supplied literature-source link labelled ``restricted geometric Langlands formulation'': \url{https://arxiv.org/abs/2010.01906}. The record does not supply an author, title, or statement locator for this reference.
\item Literature assessment dated exactly \texttt{2026-08-22}: the supplied curation labels the problem open and reports verification in elementary examples such as the projective line and a torus, with no general proof for reductive groups established. This is supplied dated curation derived from the OWR material, not an independently verified assessment; the time-sensitive status and cited claims require independent verification.
\item The registry labels the statement \texttt{corrected\_verified}, with the accompanying explanation that explicit formulas were dollar-delimited and the supplied mathematical content preserved. That supplied label does not establish mathematical correctness or resolve the missing hypotheses.
\end{itemize}
\end{document}
